%
%       LaTeX- Vereinfachungen / U.Nitsch 1994
%      ----------------------------------------
%        Liste aller Kommandos
%       -----------------------
% Kommando   anstelle von ...      Bedeutung
%
%  <=         \le                   kleiner gleich
%  <<         \ll                   viel kleiner
%  >=         \ge                   groesser gleich
%  >>         \gg                   viel groesser
%  ~a ...     \alpha ...            kleine griechische Buchstaben
%  ~A ...     \Alpha ...            grosse griechische Buchstaben
%  \8         \infty                unendlich
%  \apx       \approx               etwa
%  \be        \begin{equation}   
%  \bea       \begin{eqnarray}
%  \beas      \begin{eqnarray*}
%  \bs        \backslash            Backslash
%  \cd        \cdots                mittige Dots
%  \da        \downarrow            einfacher Pfeil nach unten
%  \dd        \ddots                diagonale Dots
%  \Do        \o                    kleines daenisches 'o'
%  \ee        \end{equation}
%  \eea       \end{eqnarray}
%  \eeas      \end{eqnarray*}
%  \eqv       \equiv                aequivalent
%  \es        \emptyset             leere Menge
%  \ex        \exists               es existiert
%  \f         \frac                 Bruch
%  \fa        \forall               fuer alle    
%  \fmf                             Rahmen um einfache Formel
%  \fsf                             Rahmen um mehrzeilige Formel
%  \inn       \int\nolimits         Integral nolimits
%  \inl       \int\limits           Integral limits
%  \l<        \langle               Winkelklammer links
%  \la        \leftarrow            einfacher Pfeil nach links
%  \ld        \ldots                line Dots
%  \lgla      \longleftarrow        langer Pfeil nach links
%  \lglra     \longleftrightarrow   langer Pfeil nach links und rechts
%  \lgmt      \longmapsto           langer folgt-aus-Pfeil
%  \lgra      \longrightarrow       langer Pfeil nach rechts
%  \lra       \leftrightarrow       einfacher Pfeil nach links und rechts
%  \Lra       \leftrightarrow       doppelter Pfeil nach links und rechts
%  \lt        \leadsto              fuehrt-zu-Pfeil
%  \mt        \mapsto               folgt-aus-Pfeil
%  \n         \not                  Negation
%  \o+        \oplus
%  \o-        \ominus
%  \o.        \opoint
%  \o/        \oslash
%  \oin       \oint\nolimits        Kreisintegral nolimits
%  \oil       \oint\limits          Kreisintegral limits
%  \ol        \overline             Ueberstrich
%  \ox        \otimes
%  \pn        \prod\nolimits        Produkt nolimits
%  \pl        \prod\limits          Produkt limits
%  \Pl        \l                    kleines polnisches 'l'
%  \pll       \parallel             parallel
%  \r>        \rangle               Winkelklammer rechts
%  \ra        \rightarrow           einfacher Pfeil nach rechts
%  \Ra        \rightarrow           doppelter Pfeil nach rechts
%  \s         \sqrt                 Wurzel
%  \sbs       \subset               Untermenge
%  \sbse      \subseteq             Untermenge und gleich
%  \sps       \supset               Obermenge
%  \spse      \supseteq             Obermenge und gleich
%  \sun       \sum\nolimits         Summe nolimits
%  \sul       \sum\limits           Summe limits
%  \tri       \triangle             Dreieck
%  \ua        \uparrow              einfacher Pfeil nach unten
%  \uda       \updownarrow          einfacher Pfeil nach unten und oben
%  \ul        \underline            Unterstrich
%  \vd        \vdots                vertikale Dots
%  \x         \times
%
%        Noetige Umdefinitionen
%       ------------------------
% Rettung der Symbole von \o und \l
%
\newsavebox{\boxA}
\sbox{\boxA}{\o}
\def\Do{\usebox{\boxA}}
\newsavebox{\boxB}
\sbox{\boxB}{\l}
\def\Pl{\usebox{\boxB}}
%
%        Mathematische Umgebungen
%       --------------------------
% Math- und Mathindent-Umgebung nicht weiter beachtet
%
\def\be{\begin{equation}}
\def\ee{\end{equation}}
\def\bea{\begin{eqnarray}}
\def\eea{\end{eqnarray}}
\def\beas{\begin{eqnarray*}}
\def\eeas{\end{eqnarray*}}
%
%        Exponenten und Indizes
%       ------------------------ 
% keine Vereinfachungen noetig
%
%        Brueche
%       ---------
\def\f{\frac}
%
%        Wurzeln
%       ---------
\def\s{\sqrt}
%
%        Auslassungen
%       --------------
\def\ld{\ldots}
\def\cd{\cdots}
\def\vd{\vdots}
\def\dd{\ddots}
%
%        Griechische Buchstaben
%       ------------------------
% Die verkuerzenden Schreibweisen der griech. Buchstaben sind aus der
% "Systematik der internat. Math.-schrift fuer Blinde" abgeleitet.
%
\catcode`\~=\active
\def~#1{\ifx#1a\alpha\fi
        \ifx#1b\beta\fi
        \ifx#1g\gamma\fi
        \ifx#1d\delta\fi
        \ifx#1e\epsilon\fi
        \ifx#1z\zeta\fi
        \ifx#1j\eta\fi
        \ifx#1h\theta\fi
        \ifx#1i\iota\fi
        \ifx#1k\kappa\fi
        \ifx#1l\lambda\fi
        \ifx#1m\mu\fi
        \ifx#1n\nu\fi
        \ifx#1x\xi\fi
        \ifx#1o o\fi
        \ifx#1p\pi\fi
        \ifx#1r\rho\fi
        \ifx#1s\sigma\fi
        \ifx#1t\tau\fi
        \ifx#1u\upsilon\fi
        \ifx#1f\phi\fi
        \ifx#1c\chi\fi
        \ifx#1y\psi\fi
        \ifx#1w\omega\fi
        \ifx#1A A\fi
        \ifx#1B B\fi
        \ifx#1G\Gamma\fi
        \ifx#1D\Delta\fi
        \ifx#1E E\fi
        \ifx#1Z Z\fi
        \ifx#1J H\fi
        \ifx#1H\Theta\fi
        \ifx#1I I\fi
        \ifx#1K K\fi
        \ifx#1L\Lambda\fi
        \ifx#1M M\fi
        \ifx#1N N\fi
        \ifx#1X\Xi\fi
        \ifx#1O O\fi
        \ifx#1P\Pi\fi
        \ifx#1R P\fi
        \ifx#1S\Sigma\fi
        \ifx#1T T\fi
        \ifx#1U\Upsilon\fi
        \ifx#1F\Phi\fi
        \ifx#1C X\fi
        \ifx#1Y\Psi\fi
        \ifx#1W\Omega\fi}
%
%        Zweistellige Operatoren
%       -------------------------
% es wird mit dem daenischen 'o' konkurriert (wurde darum umbenannt)
%
\def\o#1{\ifx#1+\oplus\fi
         \ifx#1-\ominus\fi 
         \ifx#1/\oslash\fi 
         \ifx#1.\odot\fi}
\def\ox{\otimes}
%
%        Relationen
%       ------------
\newsavebox{\boxC}
\sbox{\boxC}{$<$}
\catcode`\<=\active
\def<#1{\ifx#1=\le
        \else\ifx#1<\ll
        \else\usebox{\boxC}#1\fi\fi}
\newsavebox{\boxD}
\sbox{\boxD}{$>$}
\catcode`\>=\active
\def>#1{\ifx#1=\ge
        \else\ifx#1>\gg
        \else\usebox{\boxD}#1\fi\fi}
\def\n{\not}
\def\sbs{\subset}
\def\sps{\supset}
\def\sbse{\subseteq}
\def\spse{\supseteq}
\def\tri{\triangle}
\def\apx{\approx}
\def\eqv{\equiv}
\def\pll{\parallel}
%
%        Pfeil- oder Zeigersymbole
%       ---------------------------
% die wichtigsten Symbole wurden verkuerzt
%
\def\la{\leftarrow}
\def\ra{\rightarrow}
\def\Ra{\Rightarrow}
\def\ua{\uparrow}
\def\da{\downarrow}
\def\lra{\leftrightarrow}
\def\Lra{\leftrightarrow}
\def\uda{\updownarrow}
\def\lgla{\longleftarrow}
\def\lgra{\longrightarrow}
\def\lglra{\longleftrightarrow}
\def\mt{\mapsto}
\def\lgmt{\longmapsto}
\def\lt{\leadsto}
%
%        Vermischte Symbole
%       --------------------
\def\8{\infty}
\def\es{\emptyset}
\def\fa{\forall}
\def\ex{\exists}
\def\x{\times}
%
%        Symbole mit unterschiedlichen Groessen und Grenzangaben
%       ---------------------------------------------------------
% viele moegliche Abkuerzungen sind schon in Gebrauch
%
\def\inl{\int\limits}
\def\inn{\int\nolimits}
\def\oil{\oint\limits}
\def\oin{\oint\nolimits}
\def\pl{\prod\limits}
\def\pn{\prod\nolimits}
\def\sul{\sum\limits}
\def\sun{\sum\nolimits}
%
%        Funktionen
%       ------------
% keine Aenderungen empfehlenswert
%
%        Begrenzer
%       -----------
%
\def\l<{\langle}
\def\r>{\rangle}
\def\bs{\backslash}
%
%        Diakritische Zeichen
%       ----------------------
% Akzente, zu selten benoetigt
%
\def\ol{\overline}
\def\ul{\underline}
%
%        Gestaltungsmittel
%       -------------------
% Rahmen um Formeln, Determinanten, Matrizen
%
\def\fsf#1{\fbox{$#1$}}
\def\fmf#1{\fbox{$\begin{array}{rcl}#1\end{array}$}}
%   
\def\dtm#1#2{\left|\begin{array}{#1}#2\end{array}\right|}
\def\mat#1{\left(\begin{pmatrix}#1\end{pmatrix}\right)}
\def\kla#1#2{\left\{\begin{array}{#1}#2\end{array}\right.}
%
%        Anmerkung
%       -----------
% Die in den Saveboxen abgelegten Symbole werden auch nach 
% Schriftgroessenaenderungen in der Groesse dargestellt, die
% bei der Einspeicherung gegolten hat.
%
